\begin{longtable}{p{0.7cm}p{3.2cm}p{7.8cm}p{3.0cm}p{2.3cm}}
\caption{Definitions of the measures used in this paper (full form)}\label{tab:defs_full}\\
\toprule No. & Term & Conceptual and operational definition & Source & Applied in \\ \midrule \endfirsthead
\multicolumn{5}{l}{\footnotesize Table~\thetable\ (continued)}\\ \toprule No. & Term & Conceptual and operational definition & Source & Applied in \\ \midrule \endhead
\midrule\multicolumn{5}{r}{\footnotesize continued on next page}\\ \endfoot
\bottomrule \endlastfoot
D1 & Gross requests ($T$) & Shares a holder asked the fund to repurchase in one quarterly tender offer; the quantity the run theories call redemption requests. Operationally the number of shares ``validly tendered and not withdrawn prior to the expiration of the Offer'' reported in the final SC TO-I/A result filing, summed across share classes where listed separately. & Result filing (SC TO-I/A), Rule 13e-4; Section~\ref{sec:definitions} & Tables~\ref{tab:sample}, \ref{tab:gates}, \ref{tab:transitions}; equations~(\ref{eq:decomp})--(\ref{eq:specC}) \\ \addlinespace[2pt]
D2 & Accepted shares ($A$) & Shares the fund repurchased in the offer. Operationally the number ``accepted for purchase'' in the result filing; set equal to $T$ when the filing states that the fund accepted ``100\% of the Shares'' or purchased ``all validly tendered'' shares or paid the net asset value of all shares tendered. & Result filing (SC TO-I/A); Section~\ref{sec:definitions} & Tables~\ref{tab:gates}, \ref{tab:transitions}; equation~(\ref{eq:decomp}) \\ \addlinespace[2pt]
D3 & Rationed offer & An offer in which the fund repurchased fewer shares than were tendered and allocated pro rata. Operationally $A<T$ in the result filing (author rule); oversubscription $T/A$. Offers enlarged above the original cap and then filled in full are not rationed. & Author rule on D1--D2 & Tables~\ref{tab:gates}, \ref{tab:transitions}, \ref{tab:h3} \\ \addlinespace[2pt]
D4 & Request rate & Requests as a share of the shares eligible for the cap. Operationally $T$ divided by the share count the fund used to set its cap (cap in shares $\div$ cap in percent, from the offer's own filing); for Blue Owl's 2023--25 offers the dollar cap $\div$ the net offering price in the same document $\div$ 5\%; for Blue Owl's 2026 offers the percentage stated in the result filing (OTIC March 2026: shares outstanding net of the January tender). & Offer and result filings; author rule, Section~\ref{sec:definitions} & Tables~\ref{tab:sample}, \ref{tab:gates}; Figure~\ref{fig:requests} \\ \addlinespace[2pt]
D5 & Unfilled shares ($U_n$) & The part of a window's requests that lapsed unfilled and may be re-tendered. Operationally $U_n=T_n-A_n$ from the result filing (equation~\ref{eq:decomp}). & D1--D2; equation~(\ref{eq:decomp}) & Table~\ref{tab:transitions}; Section~\ref{sec:stockflow} \\ \addlinespace[2pt]
D6 & Lower bound on requests not carried forward ($L_{n+1}$; ``the floor'') & A bound, not an estimate, on the requests in window $n+1$ that were not carried from window $n$. Operationally $L_{n+1}=\max\{0,T_{n+1}-U_n\}$ (equation~\ref{eq:decomp}); valid for any resubmission rate; a bound on first-time requests only if no unfilled request from an earlier window re-enters after skipping one; it does not bound over-requesting. & Author-developed; equation~(\ref{eq:decomp}); Appendix~\ref{app:model}, equation~(\ref{eq:stockflow}) & Table~\ref{tab:transitions}; Sections~\ref{sec:stockflow}, \ref{sec:policy} \\ \addlinespace[2pt]
D7 & Resubmission ceiling share & The upper bound on resubmissions in a later window, as a share of that window's gross requests. Operationally $\min\{1,U_n/T_{n+1}\}$ from D1, D2 (author rule). & Author-developed; Section~\ref{sec:definitions} & Table~\ref{tab:transitions}; Section~\ref{sec:stockflow} \\ \addlinespace[2pt]
D8 & Cycle (scheduled-expiry) dating & The convention that places each offer in time by when requests had to be in. Operationally the quarter of the scheduled expiry date in the offer document; Blue Owl Technology Income's offer scheduled for 31 December 2025 and extended to 8 January 2026 is a 2025-cycle offer. & Author rule; Section~\ref{sec:design} & All tables dated by window; Figure~\ref{fig:requests} \\ \addlinespace[2pt]
D9 & Exposure to the first proration disclosure & Whether investors deciding in a window could know that a fund had rationed. Operationally, relative to Ares's amendment of 24 March 2026 (the first public proration): Straddle = window open on that date; Post = window opened after it; the 2026-cycle indicator approximates exposure dated from the first above-cap disclosure (\FirstAboveCapDate). & Author rule; Section~\ref{sec:design}; equations~(\ref{eq:specA})--(\ref{eq:specC}) & Table~\ref{tab:regressions} \\ \addlinespace[2pt]
D10 & Headroom; low headroom & A predetermined proxy for how close a fund already was to its cap before 2026. Operationally the five-percent cap minus the request rate (D4) of the fund's last 2025-cycle offer; the low-headroom group is below the median; it proxies expected rationing exposure and does not measure beliefs. & Author-developed; Section~\ref{sec:horserace} & Table~\ref{tab:regressions}; equation~(\ref{eq:specC}) \\ \addlinespace[2pt]
D11 & Listed-twin discount & The market's valuation of an exchange-listed claim on largely the same loans relative to the sponsor's mark. Operationally price (first close on or after the filing date) $\div$ NAV per share (US-GAAP \texttt{NetAssetValuePerShare}, period end) $-1$, in percent; the value known at an offer's commencement is the latest filing before the window opened. & XBRL companyfacts (SEC); Yahoo Finance closes; Section~\ref{sec:data} & Table~\ref{tab:discounts}; Figure~\ref{fig:twins}; equation~(\ref{eq:specA}) \\ \addlinespace[2pt]
D12 & Matched group and mark & The unit on which a non-traded fund's and its listed twin's valuations of the same loans are compared. Operationally the pooled schedule-of-investments lines of one vehicle sharing issuer key, spread, interest rate, lien seniority and facility type (dollar-denominated only); mark $m_g=\sum\mathrm{FV}/\sum\mathrm{Prin}$; difference in basis points of principal, weighted by the fund's fair value. & Inline-XBRL schedules of investments (US-GAAP \texttt{InvestmentOwned*} elements); author rule, Appendix~\ref{app:soi} & Table~\ref{tab:marks}; Tables~\ref{tab:marks_quarterly}, \ref{tab:marks_largest}; equation~(\ref{eq:mark}) \\ \addlinespace[2pt]
D13 & Non-accrual share & The share of a vehicle's investments on which it has stopped accruing income. Operationally the figure the filing itself states, on the filer's own denominator and basis (fair value preferred; cost or a count where that is all the filing gives), bound to the report date by the nearest ``as of'' clause. & 10-Q/10-K text; Appendix~\ref{app:soi} & Table~\ref{tab:marks}; Section~\ref{sec:marks} \\ \addlinespace[2pt]
D14 & Balance-sheet flows relative to net assets & What the fund paid out, took in and invested in a quarter, scaled by its size. Operationally US-GAAP cash-flow elements (e.g.\ \texttt{PaymentsForRepurchaseOfCommonStock}, \texttt{PaymentsToAcquireInvestments}) from 10-Q/10-K XBRL, quarterly values as differences of year-to-date amounts, divided by net assets; dispositions and repayments are reported together by the filers. & XBRL companyfacts (SEC); Section~\ref{sec:data} & Table~\ref{tab:balance}; Section~\ref{sec:balance} \\ \addlinespace[2pt]
\end{longtable}
